%% abtex2-modelo-trabalho-academico.tex, v-1.7.1 laurocesar
%% Copyright 2012-2013 by abnTeX2 group at http://abntex2.googlecode.com/
%%
%% This work may be distributed and/or modified under the
%% conditions of the LaTeX Project Public License, either version 1.3
%% of this license or (at your option) any later version.
%% The latest version of this license is in
%%   http://www.latex-project.org/lppl.txt
%% and version 1.3 or later is part of all distributions of LaTeX
%% version 2005/12/01 or later.
%%
%% This work has the LPPL maintenance status `maintained'.
%%
%% The Current Maintainer of this work is the abnTeX2 team, led
%% by Lauro C\'{e}sar Araujo. Further information are available on
%% http://abntex2.googlecode.com/
%%
%% This work consists of the files abntex2-modelo-trabalho-academico.tex,
%% abntex2-modelo-include-comandos and abntex2-modelo-references.bib
%%

% ------------------------------------------------------------------------
% ------------------------------------------------------------------------
% abnTeX2: Modelo de Trabalho Academico (tese de doutorado, dissertacao de
% mestrado e trabalhos monograficos em geral) em conformidade com
% ABNT NBR 14724:2011: Informacao e documentacao - Trabalhos academicos -
% Apresentacao
% ------------------------------------------------------------------------
% ------------------------------------------------------------------------

\input{preambulo}

% ---- compila o \'{\i}ndice  ----
\makeindex
%\makeglossaries
% ---
\usepackage[utf8]{inputenc}

% ---- In\'{\i}cio do documento ----
\begin{document}

% Retira espa\c{c}o extra obsoleto entre as frases.
\frenchspacing

% ---- ELEMENTOS PR\'{E}-TEXTUAIS ----
%\pretextual

% --- Capa ---
\imprimircapa
% ---

% --- Folha de rosto (o * indica que haver\'{a} a ficha catalogr\'{a}fica) ---
\imprimirfolhaderosto
% ---

% --- Inserir a ficha catalogr\'{a}fica ---
%\includepdf{ficha_catalografica.pdf}

% --- Para a vers\~{a}o final, delete as linhas abaixo e insira a linha do \includepdf
 \begin{fichacatalografica}
    \vspace*{\fill}
    \begin{center}
        \textsc{Inclua aqui o pdf com a ficha catalogr\'{a}fica fornecida pela BAE.}
    \end{center}
    \vspace*{\fill}
% --- --- ---
    %\includepdf{ficha-catalografica.pdf}
 \end{fichacatalografica}
% ---

% ---

% --- Inserir folha de aprova\c{c}\~{a}o ---
%\includepdf{folha_de_assinaturas_oficial.pdf}

% Isto \'{e} um exemplo de Folha de aprova\c{c}\~{a}o, elemento obrigat\'{o}rio da NBR
% 14724/2011 (se\c{c}\~{a}o 4.2.1.3). Voc\^{e} pode utilizar este modelo at\'{e} a aprova\c{c}\~{a}o
% do trabalho. Ap\'{o}s isso, substitua todo o conte\'{u}do deste arquivo por uma
% imagem da p\'{a}gina assinada pela banca com o comando abaixo:

% --- Na vers\~{a}o final, exclua essas linhas e insira o \includepdf
\newpage
\vspace*{\fill}
\begin{center}
    \textsc{Inclua aqui a folha de assinaturas.}
\end{center}
\vspace*{\fill}
\newpage
% --- --- ---
%\includepdf[pagecommand={\thispagestyle{plain}}]{folha-assinaturas.pdf}	
\cleardoublepage

% ---

%% --- Dedicat\'{o}ria ---
\begin{dedicatoria}
    \vspace*{\fill}
    \centering
    \noindent
    \textit{Dedico esta tese \`{a} todo mundo.}
    \vspace*{\fill}
\end{dedicatoria}
%% ---

%% --- Agradecimentos ---
\begin{agradecimentos}
	Escreva seus agradecimentos.
    
    
	\lipsum[1-3]
\end{agradecimentos}

%% --- RESUMOS (em portugu\^{e}s e ingl\^{e}s --- 150 a 500 palavras

\begin{resumo}
    
Insira seu resumo.

	\lipsum[1]
    
    \vspace{\onelineskip}

    \noindent\textbf{Palavras-chaves}: palavra-chave 1; palavra-chave 2; palavra-chave 3.

    \vspace{\onelineskip}
    \vspace{\onelineskip}

    \begin{otherlanguage*}{english}
    \begin{center}{\ABNTEXchapterfont\huge Abstract}\end{center}

Same content of "Resumo".       

    \vspace{\onelineskip}

    \noindent\textbf{Keywords}: keyword 1; keyword 2; keyword 3. 

    \end{otherlanguage*}
\end{resumo}
%% ---



%% ---
%
%% --- Ep\'{\i}grafe  ---
\begin{epigrafe}
    \vspace*{\fill}
	\begin{flushright}
		\textit{``Escreva aqui a sua ep\'{\i}grafe''\\
		(Cita\c{c}\~{a}o)}
	\end{flushright}
\end{epigrafe}
%% ---
%
%
%% --- inserir lista de ilustra\c{c}\~{o}es ---
\pdfbookmark[0]{\listfigurename}{lof}
\listoffigures*
\cleardoublepage
%% ---
%
%% --- inserir lista de tabelas ---
\pdfbookmark[0]{\listtablename}{lot}
\listoftables*
\cleardoublepage
%% ---
%
%% --- inserir lista de Acronimos e Abrevia\c{c}\~{o}es ---
%\renewcommand{\nomname}{Lista de Acr\^{o}nimos e Abrevia\c{c}\~{o}es}
%\pdfbookmark[0]{\nomname}{las}
%\printnomenclature
%\cleardoublepage
%% ---

%% --- inserir o sumario ---
\pdfbookmark[0]{\contentsname}{toc}
\tableofcontents*
\clearpage
%% ---

\pagenumbering{arabic}
\setcounter{page}{14}

% ---- ELEMENTOS TEXTUAIS ----
\textual

\include{introducao}
\chapter{Introdu\c{c}\~{a}o}
\label{cap:intro}
\lipsum[1-4]    % Dummy text para a introdu\c{c}\~{a}o

\part{Fundamentos (Opcional)}

\include{cap1}

\include{cap5}

%\include{capitulo7}

% --- Finaliza a parte no bookmark do PDF, para que se inicie o bookmark na raiz ---
\bookmarksetup{startatroot}%
% ---

% --- Conclus\~{a}o ---
\include{conclusao}

% ---- ELEMENTOS P\'{O}S-TEXTUAIS ----
\postextual

% ---- Refer\^{e}ncias bibliogr\'{a}ficas ----
\bibliography{tese}

% ---- Ap\^{e}ndices ----

% ---
% Inicia os ap\^{e}ndices
% ---
\begin{apendicesenv}

\include{apendice}

\end{apendicesenv}
% ---

% ---- Anexos ----

% ---
% Inicia os anexos - opcional
% ---
%\begin{anexosenv}
%
%% Imprime uma p\'{a}gina indicando o in\'{\i}cio dos anexos
%\partanexos
%
%% ---
%\chapter{Morbi ultrices rutrum lorem.}
%% ---
%\lipsum[30]
%
%\end{anexosenv}
%
%% ---- INDICE REMISSIVO ----
%
%\printindex

\end{document} 